% Figure 1.3: the four shapes a solution set can take.
\begin{tikzpicture}[scale=0.78,>=stealth,
  ax/.style={->,gray}]
  % (a) unique solution: a point
  \begin{scope}
    \draw[ax] (0,0) -- (1.7,0);\draw[ax] (0,0) -- (0,1.5);\draw[ax] (0,0) -- (-0.9,-0.7);
    \filldraw[accent] (0.55,0.6) circle (2.2pt);
    \node at (0.4,-1.35) {\small (a) point};
  \end{scope}
  % (b) one free variable: a line
  \begin{scope}[xshift=3.6cm]
    \draw[ax] (0,0) -- (1.7,0);\draw[ax] (0,0) -- (0,1.5);\draw[ax] (0,0) -- (-0.9,-0.7);
    \draw[thick] (-0.6,-0.1) -- (1.5,1.2);
    \node at (0.45,-1.35) {\small (b) line};
  \end{scope}
  % (c) two free variables: a plane
  \begin{scope}[xshift=7.4cm]
    \draw[ax] (0,0) -- (1.7,0);\draw[ax] (0,0) -- (0,1.5);\draw[ax] (0,0) -- (-0.9,-0.7);
    \draw[thick,fill=accentsoft] (-0.6,0.15) -- (0.7,0.65) -- (1.6,1.25) -- (0.3,0.75) -- cycle;
    \node at (0.5,-1.35) {\small (c) plane};
  \end{scope}
  % (d) inconsistent: parallel planes, empty set
  \begin{scope}[xshift=11.2cm]
    \draw[thick,fill=accentsoft] (-0.6,0.1) -- (0.7,0.6) -- (1.6,1.2) -- (0.3,0.7) -- cycle;
    \draw[thick,fill=graybg] (-0.5,-0.6) -- (0.8,-0.1) -- (1.7,0.5) -- (0.4,0.0) -- cycle;
    \node at (0.55,-1.35) {\small (d) empty set};
  \end{scope}
\end{tikzpicture}
